\section{Криволинейное движение}
\subsection{Равномерное движение по окружности}

%%%%%%%%%%%%%%%%%%%%%%%%%%%%%%%%%%%%%%%%%%%%%%%%%%%%
%%%%%%%%%%%%%%%%%%%%%%%%%%%%%%%%%%%%%%%%%%%%%%%%%%%%
 
 \z	{На часах 16:00. Через какое время после этого часовая и минутная стрелки часов встретятся во второй раз?}
%
{$t=\frac{8\pi}{3(\omega_\text{мин}-\omega_\text{час})}\approx 87.3\,\text{мин}$}

%%%%%%%%%%%%%%%%%%%%%%%%%%%%%%%%%%%%%%%%%%%%%%%%%%%%

 \z[{\begin{wrapfigure}[6]{r}{0.2\textwidth}%
		\vspace{-15pt}%
		\begin{flushright}%
 			\includegraphics[width=0.2\textwidth]{3.1.2}%
		\end{flushright}%
	\end{wrapfigure}%
	}]%
	{	Маленький шарик влетает со скоростью $v$ в малое отверстие в стенке полого цилиндра, вращающегося вокруг своей оси (см. рисунок). Радиус $R$ цилиндра много больше толщины его стенок. Скорость шарика перпендикулярна оси цилиндра. Какой должна быть угловая скорость вращения цилиндра $\omega$ для того, чтобы шарик вылетел наружу, не испытав соударений? Силу тяжести не учитывайте.}
%
{$\omega=\frac{(2n+1)\pi\nu}{2R}, \, n \in \mathbb{Z} $}
%%%%%%%%%%%%%%%%%%%%%%%%%%%%%%%%%%%%%%%%%%%%%%%%%%%%
\z{Точка движется по окружности с постоянной скоростью $v = 0.5 \, \text{м/с}$. Вектор скорости изменяет направление на $\Delta \varphi = 30\degree$ за время $\Delta t = 2\, \text{с}$. Каково нормальное ускорение точки $a_n$?}
%
{$a_n=\frac{v\pi}{6\Delta t}=0.13\, \text{м}/\text{с}^2$}
%%%%%%%%%%%%%%%%%%%%%%%%%%%%%%%%%%%%%%%%%%%%%%%%%%%%	
\z{С какой скоростью $v$ и в какое время суток должен лететь самолет на широте Санкт-Петербурга ($\varphi = 60\degree$), чтобы летчик видел Солнце все время на юге?}
%
{$v\approx840\, \text{км/ч}$}
%%%%%%%%%%%%%%%%%%%%%%%%%%%%%%%%%%%%%%%%%%%%%%%%%%%%
 \z{С какой скоростью должен лететь спутник, чтобы, все время «падая» на Землю с ускорением $g$, двигаться по окружности? Принять радиус орбиты $R = 6400 \,\text{км}$.}
 %
{$v=\sqrt{gR}=8\,\text{км/с}$}
